
\begin{frame}{Exemplo: sistema de controle de temperatura}
    \centering
    \resizebox{0.9\linewidth}{!}{%
    \begin{tikzpicture}[
        bloco/.style={draw, rounded corners, minimum width=2.6cm,
                      minimum height=1cm, align=center, fill=blue!10},
        seta/.style={-{Latex[length=3mm]}, thick}
    ]
        \node[bloco] (sensor) {Sensor\\Temperatura};
        \node[bloco, right=0.8cm of sensor] (adc) {ADC};
        \node[bloco, right=0.8cm of adc] (mcu) {Controlador\\digital};
        \node[bloco, right=0.8cm of mcu] (dac) {DAC ou\\PWM};
        \node[bloco, below=1.3cm of dac] (atuador) {Ventilador\\ou aquecedor};

        \draw[seta] (sensor) -- (adc);
        \draw[seta] (adc) -- (mcu);
        \draw[seta] (mcu) -- (dac);
        \draw[seta] (dac) -- (atuador);
    \end{tikzpicture}%
    }

    \vspace{0.4cm}

    O ADC permite \textbf{medir} a temperatura. O DAC ou PWM permite
    \textbf{controlar} um atuador.
\end{frame}

